(1) 对自旋1/2，\(|+x\rangle=(|+z\rangle+|-z\rangle)/\sqrt2\)。因此初态经 \(+x\) 选择的概率
\[
P(+x|+z)=|\langle+x|+z\rangle|^2=\frac12.
\]
选中后态坍缩为 \(|+x\rangle\)。

(2) 两个方向间夹角为 \(\alpha=\arccos(\hat x\cdot\boldsymbol n)=\arccos(\sin\theta)\)。自旋方向投影概率为 \((1+\cos\alpha)/2\)，故
\[
P(+n|+x)=\frac{1+\sin\theta}{2}.
\]

(3) 第二次选中后的态为 \(|+n\rangle\)，在 z 方向再次测得正自旋的概率 \(P(+z|+n)=(1+\cos\theta)/2\)。因此三次均通过指定通道的联合概率为
\[
P_{\rm joint}=\frac12\frac{1+\sin\theta}{2}\frac{1+\cos\theta}{2}
=\frac{(1+\sin\theta)(1+\cos\theta)}8.
\]
其中每一步是上一步选择后的条件概率，不能把角度直接加在同一个投影公式里。

(4) 若没有中间的 x、n 测量，初始 \(|+z\rangle\) 直接测得 \(+z\) 的概率为1；加入两次选择会降低全程通过的比例并改变量子态。非对易投影测量通常不可交换，测量过程不只是读取某个预先存在的经典三维自旋矢量。
